\subsection{连续作用于局部紧空间的群}

命题 9. 设 G 是连续作用于局部紧空间 X 的拓扑群。那么，若 X/G 是豪斯多夫的，则它是局部紧的。

由于 G 在 X 上定义的等价关系是开的（§ 2，第 4 节，引理 2），本命题由第一章 § 10 第 4 节命题 10 得出。

命题 10. 设 G 是连续作用于局部紧空间 X 的拓扑群，且设 X/G 是豪斯多夫的。设 \(\varphi\) 是 X 到 X/G 上的典范映射。那么，若 K' 是 X/G 的任意紧子集，则存在 X 的紧子集 K 使得 \(\varphi(K) = K'\)。

由于由 G 定义的等价关系是开的（§ 2，第 4 节，引理 2），本命题是第一章 § 10 第 4 节命题 10 的特殊情形。

命题 11. 设 G 是逆紧作用于非空空间 X 的豪斯多夫拓扑群。若 X 是紧的（相应地，局部紧的），则 G 与 X/G 也是紧的（相应地，局部紧的）。

由假设，映射 \(\theta : (s, x) \to (x, s.x)\)（\(G \times X\) 到 \(X \times X\) 中的映射）是逆紧的；若 \(X \times X\) 是紧的（相应地，局部紧的），则第一章 § 10 第 4 节命题 9 的推论表明 \(G \times X\) 也是紧的（相应地，局部紧的），从而 \(G\) 也是紧的（相应地，局部紧的），因为 \(X \neq \emptyset\)。由于 \(X/G\) 是豪斯多夫的（第 2 节，命题 3），\(X\) 的紧性（相应地，局部紧性）蕴含 \(X/G\) 的紧性（相应地，局部紧性）{[}第一章，§ 10，第 4 节，命题 8（相应地，命题 9）{]}（见 § 2，习题 29）。

我们现在给出一些判别法，使我们能够断言豪斯多夫拓扑群 G 逆紧作用于局部紧空间 X。对 X 的每一对子集 K, L，我们用 P(K, L) 表示由所有这样的 \(s \in G\) 组成的集合：它们满足 \(s \cdot K \cap L \neq \emptyset\)。

定理 1. 设 G 是连续作用于拓扑空间 X 的豪斯多夫拓扑群。设 K 是 X 的紧子集，L 是 X 的闭子集。那么：

\begin{enumerate}
\def\labelenumi{\alph{enumi})}
\item
  集合 \(\mathbf{P}(\mathbf{K},\mathbf{L})\) 在 \(\mathbf{G}\) 中是闭的。
\item
  若 G 逆紧作用于 X 且 L 是紧的，则 P(K, L) 是紧的。
\item
  反之，若 X 是局部紧的，并且对 X 的每一对紧子集 K, L，P(K, L) 都在 G 中相对紧{[}从而由 a) 是紧的{]}；则 G 逆紧作用于 X（且若 X 非空，则由命题 11，G 是局部紧的）。
\end{enumerate}

映射 \((s, x) \to s.x\)（\(G \times K\) 到 X 中的映射）是连续的，因此 L 在此映射下的逆像 \(L'\) 是闭的。由于 K 是紧的，投影 \(pr_{1}: G \times K \to G\) 是逆紧的（第一章，§ 10，第 2 节，定理 1 的推论 5），因此 \(L'\) 在 \(pr_{1}\) 下的像是闭的。这个像就是 \(P(K, L)\)，于是 a) 得证。

为证 b)：X 是豪斯多夫的（第 2 节，命题 3）。由假设，映射 \(\theta: (s, x) \to (x, s.x)\)（\(G \times X\) 到 \(X \times X\) 中的映射）是逆紧的；由于 \(K \times L\) 是紧的，\(\theta^{-1}(K \times L)\) 也是紧的，因为 \(X\) 是豪斯多夫的（第一章，§ 10，第 2 节，命题 6）。于是 \(P(K, L)\)（\(\theta^{-1}(K \times L)\) 在 \(G\) 中的投影）是紧集。

为证 c)：由于 \(K \times L\) 在 \(X \times X\) 中闭，由此推出 \(\theta^{-1}(K \times L)\) 在 \(P(K, L) \times K\) 中闭，从而在 c) 的假设下是紧的。由于 \(X \times X\) 的每个紧子集都含于形如 \(K \times L\) 的紧子集中，由此推出 \(X \times X\) 的任意紧子集在 \(\theta\) 下的逆像是紧的；又由于 \(X \times X\) 是局部紧的，这就表明 \(\theta\) 是逆紧的（第一章，§ 10，第 3 节，命题 7）{[}见第四章，§ 1，习题 4c){]}。

注. 显然我们有 \(\mathrm{P}(\mathrm{K},\mathrm{L})\subset\mathrm{P}(\mathrm{K}\cup\mathrm{L},\mathrm{K}\cup\mathrm{L})\)；因此，为使 G 逆紧作用于局部紧空间 X，只需对 X 的每个紧子集 K，集合 \(\mathrm{P}(\mathrm{K},\mathrm{K})\) 在 G 中相对紧即可。特别地，离散群 G 逆紧作用于局部紧空间 X 当且仅当：对 X 的每个紧子集 K，由所有这样的 \(s\in G\) 组成的集合是有限的：它们满足 \(s.K\cap K\neq\emptyset\)。

例. * 设 X 是一个复解析流形，解析同构于 \(\mathbf{C}^n\) 的一个有界开子集，并设 G 是 X 的解析自同构群。紧收敛拓扑与 G 的群结构相容，并且可以证明 G 逆紧作用于 X。特别地，G 的每个离散子群都逆紧作用于 X。

例如取 X 为上半平面 \(\mathfrak{J}(z)>0\)，它解析同构于 C 中的一个开圆盘。此时 G 是由所有形如 \(z\to (az+b)/(cz+d)\) 的变换组成的群，其中 \(a,b,c,d\) 是实数且 \(ad-bc\neq0\)。G 中由所有满足 \(a,b,c,d\) 为整数且 \(ad-bc=1\) 的此类变换组成的子群 H 是 G 的离散子群，称为模群。根据上文所述，它逆紧作用于上半平面 \(\mathfrak{J}(z)>0\)。 *

命题 12. 设 G 是连续作用于拓扑空间 X 的豪斯多夫拓扑群。设 K 是 X 的紧子集，并设 \(\rho_{K}\) 是映射 \((s, x) \to s.x\)（\(G \times K\) 到 X 中的映射）。那么：

\begin{enumerate}
\def\labelenumi{\alph{enumi})}
\item
  若 G 逆紧作用于 X，则 \(\rho_{\mathbf{K}}\) 是逆紧的。
\item
  若 X 是局部紧的并且对 X 的每个紧子集 K，\(\rho_{K}\) 都是逆紧的，则 G 逆紧作用于 X。
\end{enumerate}

映射 \(\rho_{\mathbf{K}}\) 可分解为 \(\mathbf{G} \times \mathbf{K} \xrightarrow{\theta_{\mathbf{K}}} \mathbf{K} \times \mathbf{X} \xrightarrow{\mathrm{pr}_2} \mathbf{X}\)，其中 \(\theta_{\mathbf{K}}\) 是 \(\mathbf{G} \times \mathbf{K}\) 上的这样一个限制：被限制的映射是 \(\theta: (s, x) \to (x, s.x)\)（\(\mathbf{G} \times \mathbf{X}\) 到 \(\mathbf{X} \times \mathbf{X}\) 中的映射）。由于 \(\theta^{-1}(\mathbf{K} \times \mathbf{X}) = \mathbf{G} \times \mathbf{K}\)，\(\theta_{\mathbf{K}}\) 在 \(\theta\) 逆紧时是逆紧的（第一章，§ 10，第 1 节，命题 3）。另一方面，由于 \(\mathbf{K}\) 是紧的，投影 \(\mathrm{pr}_2: \mathbf{K} \times \mathbf{X} \to \mathbf{X}\) 是逆紧的（第一章，§ 10，第 2 节，定理 1，推论 5），于是 \(\rho_{\mathbf{K}}\) 是逆紧的（同处，第 1 节，命题 5）。

反之，设 \(\rho_{\mathbf{K}}\) 对每个紧子集 \(\mathbf{K} \subset \mathbf{X}\) 都是逆紧的。若 \(\mathbf{L}\) 是 \(\mathbf{X}\) 的紧子集，则 \(\rho_{\mathbf{K}}^{-1}(\mathbf{L})\) 是 \(\mathbf{G} \times \mathbf{K}\) 的紧子集，它在 \(\mathbf{G}\) 中的投影是 \(\mathrm{P}(\mathbf{K}, \mathbf{L})\)；因此 \(\mathrm{P}(\mathbf{K}, \mathbf{L})\) 是紧的。于是若 \(\mathbf{X}\) 是局部紧的，则由定理 1 推出 \(\mathbf{G}\) 逆紧作用于 \(\mathbf{X}\)。

推论. 设 G 是逆紧作用于拓扑空间 X 的豪斯多夫拓扑群。若 K 是 X 的任意紧子集且 F 是 G 的任意闭子集，则 F.K 在 X 中闭。

这由命题 12 以及第一章 § 10 第 1 节命题 1 得出。

